% ---------------------------------------------------------------------------
% The scattering-strength expansion, re-derived.
%
% SELF-CONTAINED BY DESIGN. This file inputs one generated table and defines
% nothing; it can be moved into Paper 2 by moving two files, this and
% tables/scattering-expansion-v1.3.0.tex. Decision S42.
% ---------------------------------------------------------------------------
\section{The scattering-strength expansion, re-derived}
\label{sec:scattering-expansion}

\emph{Annals~II} Appendix~C writes the almost-Friedmann--Lema\^{\i}tre
integrated Boltzmann equation for Thomson scattering as
\begin{equation}
  \mathcal B+\dot\tau+e^a\mathrm D_a\tau \;\simeq\; t_c^{-1}\big[v^ae_a-\tau\big],
  \label{eq:ibe-thomson}
\end{equation}
inverts it, and approximates the result \emph{``systematically \dots in terms
of the smallness parameter $t_c$''} by the recursion
\begin{equation}
  \tau_{(n)} \;\simeq\; v^ae_a-t_c\Big[\mathcal B+\dot\tau_{(n-1)}
  +e^a\mathrm D_a\tau_{(n-1)}\Big],
  \label{eq:appc-recursion}
\end{equation}
printing $\tau_{(0)}$, $\tau_{(1)}$ and $\tau_{(2)}$. Section~5.5.3,
\emph{Expansion in Thompson Scattering Time}, states that the scheme carries
\emph{``terms up to $\mathcal O(t_c^3)$''}, so the ceiling is the antecedent's
and the printed orders stop one short of it.

\subsection{The recursion is geometric}

Write $\mathcal D\equiv\partial_t+e^a\mathrm D_a$ for the derivative along the
photon path --- which is exactly the pair of terms \eqref{eq:appc-recursion}
applies to $\tau_{(n-1)}$ --- and $V\equiv v^ae_a$. Then
\eqref{eq:ibe-thomson} reads $\tau=V-t_c(\mathcal B+\mathcal D\tau)$, so
\begin{equation}
  \big(1+t_c\mathcal D\big)\tau=V-t_c\mathcal B,
  \qquad
  \tau=\big(1+t_c\mathcal D\big)^{-1}\big(V-t_c\mathcal B\big),
  \label{eq:closed-operator}
\end{equation}
and expanding the inverse as a Neumann series in $(-t_c\mathcal D)$ gives a
closed form at every order:
\begin{equation}
  \boxed{\;\tau_{(n)}\;=\;V\;+\;\sum_{m=1}^{n}(-1)^m\,t_c^m
  \Big[\mathcal D^{m}V+\mathcal D^{m-1}\mathcal B\Big].\;}
  \label{eq:closed-form}
\end{equation}

\textbf{This shows that} Appendix~C's recursion is not an open-ended
approximation scheme but a geometric series with an explicit sum, and that
$\tau_{(1)}$ and $\tau_{(2)}$ as printed are its first two partial sums.
Equation~\eqref{eq:closed-form} is implemented in
\texttt{functions/spectra/scattering\_expansion.py} and checked against the
printed $\tau_{(1)}$ and $\tau_{(2)}$ term for term, with a deliberately
failing control, in \texttt{tests/test\_scattering\_expansion.py}. The closed
form was derived from \eqref{eq:appc-recursion} and \emph{then} compared with
the page, which is the only order in which the comparison carries information.

\subsection{Two consequences the closed form makes immediate}

\paragraph{The term counts are not chosen.} $\mathcal D^m$ expands into a sum
over words in $\{\partial_t,\;e^a\mathrm D_a\}$, so the $t_c^m$ bracket of
\eqref{eq:closed-form} holds exactly $2^m+2^{m-1}$ terms. Appendix~C prints
three at first order and six at second; those are $2+1$ and $4+2$.
Table~\ref{tab:scattering-expansion} lists every term to
$\mathcal O(t_c^3)$, generated from the same module rather than transcribed.

\paragraph{Order $m$ sources multipoles up to $\ell=m+1$.} Each
$e^a\mathrm D_a$ carries one direction vector and $V=v^ae_a$ carries one
already, so a word with $g$ gradients seeded by $V$ has $g+1$ factors of
$e^a$; reduced by Appendix~C's own solid-angle identities, which lower the rank
two at a time, its highest PSTF part has rank $g+1$. At second order that
reaches \textbf{rank three}, in the single term
$e^ce^ae^b\mathrm D_c\mathrm D_av_b$.

Equivalently, and more usefully for the high-$\ell$ question: each gradient
contributes one power of $k$, so a term of rank $r$ carries $k^{\,r-1}$. The
expansion parameter is therefore $t_ck$ --- the ratio of the mean free path to
the wavelength --- and \textbf{not} $t_c$ alone. \textbf{This shows that} the
truncation is $\ell$-dependent by construction: a cut in rank is a cut in
$t_ck$, and the order at which the expansion may be stopped falls as $\ell$
rises.

\subsection{What this costs the oscillator}

The near-tight-coupling set of \S7.1 is introduced \emph{``rewritten to first
order''} and closes with $\tau^{A_\ell}\simeq0$ for all $\ell>2$. At first
order the expansion reaches rank two, so the closure costs nothing and the
antecedent is consistent. \textbf{What the rank counting fixes is the price of
going further}: second order reaches the octupole and
$\mathcal O(t_c^3)$ the hexadecapole, and the closure carried with (152) and
(153) reaches neither.

That is approximation~\textbf{A-4} of
\texttt{provenance/SOURCE-APPROXIMATIONS-v1.md} stated as a rank rather than as
a dropped term. The term the oscillator derivation discards,
$-\tfrac25 t_c(1+R)^{-1}\mathrm D^b\mathrm D_{\langle a}u^m_{b\rangle}$, is
dropped because \emph{``the matter shear is already at least first order''} ---
a cross term of $\mathcal O(t_c\epsilon)$ between \S5.5.1's almost-Friedmann--Lema\^{\i}tre
ordering and \S5.5.3's scattering-time ordering. \emph{Annals~II}'s
Appendix~E, after Ellis--Treciokas--Matravers, states what the discarded
harmonics hold: if any four consecutive ones vanish exactly then
$\sigma_{ab}=0$ exactly. \textbf{This shows that} the truncation and the shear
are not independent --- the residual in the truncated harmonics is where the
shear lives, and a treatment that sets them to zero has set the shear to zero
with them.

\subsection{What appears at second order that does not appear at first}

Appendix~C's own solid-angle reduction gives the rank decomposition of the
source,
\begin{equation}
  \int_{4\pi}\!\mathrm d\Omega\;O^{A_\ell}\mathcal B\simeq
  \delta^\ell_0\Big[-\tfrac{4\pi}{3}\mathrm D_a\tau^a\Big]
  +\delta^\ell_1\Big[\tfrac{4\pi}{3}\big(\tfrac14\mathrm D^{a_1}\ln\rho_R+A^{a_1}\big)\Big]
  +\delta^\ell_2\Big[\tfrac{8\pi}{15}\sigma^{a_1a_2}\Big],
  \label{eq:b-ranks}
\end{equation}
so $\mathcal B$ carries the \textbf{shear} in its rank-two part. At first order
$-t_c\mathcal B$ therefore reproduces the quadrupole
$\tau_{ab}\simeq-\dot\kappa^{-1}[\sigma_{ab}+\mathrm D_{\langle a}v_{b\rangle}^B]$
of \S7.1, as it must.

At second order the closed form contributes $+t_c^2\,\mathcal D\mathcal B$, and
through \eqref{eq:b-ranks} that is $\dot\sigma_{ab}$ at rank two and
$e^c\mathrm D_c\sigma_{ab}$ at rank three. Appendix~G's
$\dot\sigma_{ab}+2H\sigma_{ab}+E_{ab}\approx0$ then gives
\begin{equation}
  \dot\sigma_{ab}=-2H\sigma_{ab}-E_{ab},
  \label{eq:sigma-dot}
\end{equation}
so \textbf{the electric Weyl tensor enters the radiation quadrupole at second
order in the scattering time}, and does not at first. \textbf{This shows that}
the second-order expansion is not merely a smaller correction of the same kind:
it introduces a gravitational tidal driving of the quadrupole that the
first-order treatment has no term for.

\paragraph{An open question, stated as a mechanism rather than a result.}
\emph{Annals~II} flags the regime itself, in a footnote to the sentence about
removing the radiation quadrupole: \emph{``If there are any anisotropic
contributions such as anisotropic scattering (in the matter frame) or large
shear (from gravitational waves) at that time this sort of approximation should
be considered with care --- such phenomena would break tight-coupling.''} For a
tensor mode both $\sigma^{(T)}_{ab}$ and $E^{(T)}_{ab}$ oscillate at the
gravitational-wave frequency, so \eqref{eq:sigma-dot} drives the second-order
quadrupole oscillatorily on a scale set by the wave and not by the sound
horizon. Whether that survives projection as a distinguishable \emph{ringing}
is open; it is recorded as question~\textbf{Q10} rather than claimed, because
nothing here computes it.

\paragraph{What this expansion cannot produce.} \eqref{eq:appc-recursion} is
\textbf{linear} in $V$ and $\mathcal B$, so \eqref{eq:closed-form} yields terms
of order $t_c^m\epsilon$ and never $\epsilon^2$. Terms of
$\mathcal O(\varepsilon^2\ell)$ require the quadratic parts of the Boltzmann
equation that \S5.5.1's almost-Friedmann--Lema\^{\i}tre linearisation discards,
and are the subject of a later release rather than of this appendix. The
high-$\ell$ sensitivity established above is a separate statement and rests on
$t_ck$, not on $\epsilon^2$.

% GENERATED by scripts/make_tables.py from config/implementation-register.toml.
% Do not edit by hand.

\begingroup
\footnotesize
\setlength{\tabcolsep}{4pt}
\renewcommand{\arraystretch}{1.25}
\begin{longtable}{llll}
\caption{\emph{Annals~II} Appendix~C's scattering-strength expansion to $\mathcal O(t_c^3)$, generated from \texttt{functions/spectra/scattering\_expansion.py}. \textbf{Rank} is the highest PSTF multipole the term can source: one per $e^a\mathrm D_a$, plus one for $V=v^ae_a$. The closure (L0082) keeps $\ell\le2$.}\label{tab:scattering-expansion}\\
\toprule
\textbf{Order} & \textbf{} & \textbf{Term} & \textbf{Rank} \\
\midrule
\endfirsthead
\toprule
\textbf{Order} & \textbf{} & \textbf{Term} & \textbf{Rank} \\
\midrule
\endhead
\bottomrule
\endfoot
$t_c^{0}$ & $+$ & $v_a e^a$ & 1 \\
$t_c^{1}$ & $-$ & $\mathcal{B}$ & 0 \\
$t_c^{1}$ & $-$ & $\big(v_a e^a\big)^{\textstyle\cdot}$ & 1 \\
$t_c^{1}$ & $-$ & $e^c \mathrm{D}_c\big(v_a e^a\big)$ & 2 \\
$t_c^{2}$ & $+$ & $\big(\mathcal{B}\big)^{\textstyle\cdot}$ & 0 \\
$t_c^{2}$ & $+$ & $e^c \mathrm{D}_c\big(\mathcal{B}\big)$ & 1 \\
$t_c^{2}$ & $+$ & $\big(\big(v_a e^a\big)^{\textstyle\cdot}\big)^{\textstyle\cdot}$ & 1 \\
$t_c^{2}$ & $+$ & $e^c \mathrm{D}_c\big(\big(v_a e^a\big)^{\textstyle\cdot}\big)$ & 2 \\
$t_c^{2}$ & $+$ & $\big(e^c \mathrm{D}_c\big(v_a e^a\big)\big)^{\textstyle\cdot}$ & 2 \\
$t_c^{2}$ & $+$ & $e^c \mathrm{D}_c\big(e^c \mathrm{D}_c\big(v_a e^a\big)\big)$ & 3 \\
$t_c^{3}$ & $-$ & $\big(\big(\mathcal{B}\big)^{\textstyle\cdot}\big)^{\textstyle\cdot}$ & 0 \\
$t_c^{3}$ & $-$ & $e^c \mathrm{D}_c\big(\big(\mathcal{B}\big)^{\textstyle\cdot}\big)$ & 1 \\
$t_c^{3}$ & $-$ & $\big(e^c \mathrm{D}_c\big(\mathcal{B}\big)\big)^{\textstyle\cdot}$ & 1 \\
$t_c^{3}$ & $-$ & $e^c \mathrm{D}_c\big(e^c \mathrm{D}_c\big(\mathcal{B}\big)\big)$ & 2 \\
$t_c^{3}$ & $-$ & $\big(\big(\big(v_a e^a\big)^{\textstyle\cdot}\big)^{\textstyle\cdot}\big)^{\textstyle\cdot}$ & 1 \\
$t_c^{3}$ & $-$ & $e^c \mathrm{D}_c\big(\big(\big(v_a e^a\big)^{\textstyle\cdot}\big)^{\textstyle\cdot}\big)$ & 2 \\
$t_c^{3}$ & $-$ & $\big(e^c \mathrm{D}_c\big(\big(v_a e^a\big)^{\textstyle\cdot}\big)\big)^{\textstyle\cdot}$ & 2 \\
$t_c^{3}$ & $-$ & $e^c \mathrm{D}_c\big(e^c \mathrm{D}_c\big(\big(v_a e^a\big)^{\textstyle\cdot}\big)\big)$ & 3 \\
$t_c^{3}$ & $-$ & $\big(\big(e^c \mathrm{D}_c\big(v_a e^a\big)\big)^{\textstyle\cdot}\big)^{\textstyle\cdot}$ & 2 \\
$t_c^{3}$ & $-$ & $e^c \mathrm{D}_c\big(\big(e^c \mathrm{D}_c\big(v_a e^a\big)\big)^{\textstyle\cdot}\big)$ & 3 \\
$t_c^{3}$ & $-$ & $\big(e^c \mathrm{D}_c\big(e^c \mathrm{D}_c\big(v_a e^a\big)\big)\big)^{\textstyle\cdot}$ & 3 \\
$t_c^{3}$ & $-$ & $e^c \mathrm{D}_c\big(e^c \mathrm{D}_c\big(e^c \mathrm{D}_c\big(v_a e^a\big)\big)\big)$ & 4 \\
\end{longtable}
\endgroup


\paragraph{One transcription note.} The fifth term of Appendix~C's $t_c^2$
bracket is printed as $e^ce^d\mathrm D_c\dot v_b$, which leaves both $d$ and
$b$ free in a scalar. It is $e^ce^d\mathrm D_c\dot v_d$ --- the time derivative
taken first and the gradient outside --- forced independently by the four words
of $\mathcal D^2$, by the bracket's count of six with the other five
unambiguous, and by the two direction vectors having nothing else to contract
with. Recorded as \textbf{B-15}; no computed quantity depends on it.
